% SEGMENT OLP-0610-S01
% Part: methods
% Chapter: proofs
% Section: reading-proofs

\documentclass[../../../include/open-logic-section]{subfiles}

\begin{document}

\olfileid{mth}{prf}{rea}
\olsection{நிறுவல்களை வாசித்தல்}

பாடநூல்களிலும் ஆய்வுக் கட்டுரைகளிலும் காணும் நிறுவல்கள்,
இதுவரை நம் எடுத்துக்காட்டுகளில் சேர்த்துள்ள விவரங்களை
அரிதாகவே முழுவதுமாகத் தருகின்றன. நிகழ்வுகளைப் பிரிக்கும்போதோ,
மறைமுக நிறுவலைக் கொடுக்கும்போதோ, ஒரு வரையறையைப்
பயன்படுத்தும்போதோ ஆசிரியர்கள் அதைத் தனியாகச் சுட்டாமல்
போகலாம். ஆகவே பாடநூலில் ஒரு நிறுவலை வாசிக்கும்போது,
அதைப் புரிந்துகொள்ள இவ்விவரங்களை நீங்களே நிரப்ப
வேண்டியிருக்கும். நிறுவலில் செய்ய வேண்டிய வெவ்வேறு
நகர்வுகளைப் பழகிக்கொள்ளவும் இது நல்ல பயிற்சி.
ஓர் எடுத்துக்காட்டைப் பார்ப்போம்.

\begin{prop}[உறிஞ்சுதல் விதி]
எல்லாக் கணங்கள் $A$, $B$ க்கும்,
\[
A \cap (A \cup B) = A
\]
\end{prop}

% SEGMENT OLP-0610-S02
\begin{proof}
$z \in A \cap (A \cup B)$ எனில் $z \in A$;
ஆகவே $A \cap (A \cup B) \subseteq A$. இப்போது
$z \in A$ என்க. அப்போது $z \in A \cup B$
உம், எனவே $z \in A \cap (A \cup B)$ உம் ஆகும்.
\end{proof}

% SEGMENT OLP-0610-S03
மேலுள்ள உறிஞ்சுதல் விதியின் நிறுவல் மிகவும் சுருக்கமானது.
பயன்படும் வரையறைகள் குறிப்பிடப்படவில்லை; ``நாம் இதை
நிறுவ வேண்டும்'' என்ற முன்னுரைகளும் இல்லை. அதை
விரித்துப் பார்ப்போம். நிறுவப்படும் முன்மொழிவு தன்னிச்சையான
கணங்கள் $A$, $B$ பற்றிய பொதுக் கூற்று. நிறுவலில் வரும்
$A$, $B$ ஆகியவை தன்னிச்சையான கணங்களைக் குறிக்கும்
மாறிகள். கணங்களின் சமத்துவத்தை நிறுவத் தேவையான
பொதுக் கூற்றுகளையே நிறுவல் தருகிறது: சமத்துவத்தின்
இடப்புறத்தின் ஒவ்வோர் !!{element} உம் வலப்புறத்தின்
!!a{element} ஆகவும், மறுதிசையிலும் இருக்க வேண்டும்.

\begin{quote}
``$z \in A \cap (A \cup B)$ எனில், $z \in A$;
  ஆகவே $A \cap (A \cup B) \subseteq A$.''
\end{quote}

இது சமத்துவ நிறுவலின் முதல் பாதி: தன்னிச்சையான~$z$
இடப்புறத்தின் !!a{element} எனில் வலப்புறத்தின்
!!a{element} என்றும், அதாவது
$A \cap (A \cup B) \subseteq A$ என்றும் காட்டுகிறது.
$z \in A \cap (A \cup B)$ என்று எடுத்துக்கொள்க.
இரு கணங்களின் வெட்டில் $z$ ஓர் !!{element} எனில்,
எனில் மட்டுமே, இரண்டிலும் அது ஓர் !!{element}.
எனவே $z \in A$ என்றும் $z \in A \cup B$
என்றும் முடிவு செய்யலாம். குறிப்பாக $z \in A$;
இதைத்தான் காட்ட விரும்பினோம். முதல் பாதிக்குத்
தேவையானது இவ்வளவுதான் என்பதால், நிறுவலின் மீதிப்
பகுதி இரண்டாம் பாதியை, அதாவது
$A \subseteq A \cap (A \cup B)$ என்பதை,
நிறுவுவதாக இருக்க வேண்டும்.

\begin{quote}
``இப்போது $z \in A$ என்க. அப்போது $z \in A \cup B$
  உம், எனவே $z \in A \cap (A \cup B)$ உம் ஆகும்.''
\end{quote}

% SEGMENT OLP-0610-S04
$z \in A$ என்று எடுத்துக்கொண்டு தொடங்குகிறோம்:
ஏனெனில் எந்த~$z$ க்கும், $z \in A$ எனில்
$z \in A \cap (A \cup B)$ என்று காட்டுகிறோம்.
$z \in A \cap (A \cup B)$ என்று காட்ட,
``$\cap$'' இன் வரையறைப்படி (i) $z \in A$
மற்றும் (ii) $z \in A \cup B$ ஆகிய இரண்டையும்
காட்ட வேண்டும். இங்கே (i) நம் எடுகோளே; அதற்கு
மேலும் நிறுவ வேண்டியதில்லை. அதனால்தான் நிறுவலில்
அது மறுபடியும் குறிப்பிடப்படவில்லை. (ii) க்கு,
கணங்களின் ஒன்றிப்பில் $z$ ஓர் !!a{element} எனில்,
எனில் மட்டுமே, அவற்றுள் குறைந்தது ஒன்றின்
!!{element} என்று நினைவுகூர்க. $z \in A$ என்பதாலும்,
$A \cup B$ என்பது $A$, $B$ ஆகியவற்றின்
ஒன்றிப்பு என்பதாலும் இங்கே அந்நிபந்தனை நிறைவேறுகிறது.
ஆகவே $z \in A \cup B$. (i) $z \in A$,
(ii) $z \in A \cup B$ இரண்டையும் காட்டிவிட்டோம்;
எனவே ``$\cap$'' இன் வரையறையின்படி
$z \in A \cap (A \cup B)$. நிறுவலில் இவ்வரையறைகள்
குறிப்பிடப்படவில்லை; வாசகர் அவற்றை ஏற்கெனவே
உள்வாங்கியிருப்பார் எனக் கருதப்படுகிறது. அப்படி
இல்லாவிட்டால் அவற்றைத் திரும்பப் பார்த்து நினைவுபடுத்திக்
கொள்ள வேண்டும். பிறகு $z \in A$ இலிருந்து
$z \in A \cup B$ எப்படி வருகிறது என்பதையும்,
$z \in A$ மற்றும் $z \in A \cup B$ ஆகியவற்றிலிருந்து
$z \in A \cap (A \cup B)$ எப்படி வருகிறது என்பதையும்
நீங்களே கண்டறிய வேண்டும்.

% SEGMENT OLP-0610-S05
எல்லாவற்றையும் வெளிப்படையாக்கிய அதே நிறுவலின்
மற்றொரு வடிவம் இதோ:
\begin{proof}{}
[கணங்களுக்கான $=$ இன் வரையறைப்படி,
$A \cap (A \cup B) = A$ என்று காட்ட,
(a) $A \cap (A \cup B) \subseteq A$ மற்றும்
% SOURCE CORRECTION TA-MTH-005: இரண்டாம் உட்கணத் திசையைத் திருத்துக.
(b) $A \subseteq A \cap (A \cup B)$ என்று
காட்ட வேண்டும். (a): $\subseteq$ இன் வரையறைப்படி,
$z \in A \cap (A \cup B)$ எனில் $z \in A$
என்று காட்ட வேண்டும்.] $z \in A \cap (A \cup B)$
எனில் $z \in A$ [$\cap$ இன் வரையறைப்படி,
$z \in A \cap (A \cup B)$ எனில், எனில் மட்டுமே,
$z \in A$ மற்றும் $z \in A \cup B$]; ஆகவே
$A \cap (A \cup B) \subseteq A$.
[(b): $\subseteq$ இன் வரையறைப்படி,
$z \in A$ எனில் $z \in A \cap (A \cup B)$
என்று காட்ட வேண்டும்.] இப்போது [(1)] $z \in A$
என்க. அப்போது [(2)] $z \in A \cup B$
[(1) இன்படி $z \in A$ அல்லது $z \in B$;
அதுவே $\cup$ இன் வரையறைப்படி $z \in A \cup B$]
உம் ஆகும். எனவே $z \in A \cap (A \cup B)$
[$\cap$ இன் வரையறைக்கு $z \in A$, அதாவது
(1), மற்றும்
% SOURCE CORRECTION TA-MTH-006: கணிதச் சின்னத் தொகுதிக்குள் கூடுதல் மூடும் அடைப்பை நீக்குக.
$z \in A \cup B$, அதாவது (2), தேவை].
\end{proof}

% SEGMENT OLP-0610-S06
\begin{prob}
$A \cup (A \cap B) = A$ என்பதற்கான பின்வரும்
நிறுவலை விரித்தெழுதுக: பயன்படுத்தப்படும் எல்லா
உய்த்தறிதல் வடிவங்களையும், ஒவ்வொரு படியும்
எந்த எடுகோளிலிருந்து அல்லது ஏற்கெனவே நிறுவிய
கூற்றிலிருந்து வருகிறது என்பதையும், எந்த இடத்தில்
எந்த வரையறையை நாட வேண்டும் என்பதையும் குறிப்பிடுக.
\begin{proof}
  $z \in A \cup (A \cap B)$ எனில் $z \in A$
  அல்லது $z \in A \cap B$. $z \in A \cap B$
  எனில் $z \in A$. எந்த $z \in A$ உம்
  $\in A \cup (A \cap B)$ ஆகும்.
\end{proof}
\end{prob}

\end{document}
